\scanpage{61}
L'ens. des val. prises par $f(z)$ dans un par. des périodes est égal à l'ensemble des valeurs prises dans tout le plan ;

\margin{$\times$}Si une fonction elliptique n'a pas de pole, elle est constante.

Remarquer que la somme, le produit, le quotient, la dérivée, de fonction elliptique de $\unsure{\overline{\mathrm m}^{\mathrm s}|=^{\mathrm s}}$ périodes ont $\overline{\mathrm m}$ périodes.

\margin{$\uparrow$}\uline{Théorème $\overline{\mathrm{II}}$} La somme des résidus d'une fonction elliptique dans un par. de périodes est nulle.

Donc, une fonction elliptique admet dans un par. de périodes au moins deux poles (chacun compté avec son ordre de multiplicité)

\uline{Théorème $\overline{\mathrm{III}}$} Le nombre des zéros et le nb des poles d'une fonction elliptique dans un parallélogramme de période est le même (en supposant tjrs que sur les bords --) ; il en est de même du nb de fois qu'il pd \deleted{la valeur} \inserted{une valeur} donnée (sous les $\overline{\mathrm m}$ réserves)

Ce nombre est appellé ordre de la fonction elliptique.

\margin{$\downarrow$}\uline{Théorème $\overline{\deleted{\unsure{\mathrm{III}|\mathrm{II}|\mathrm I}}}$} Deux fonctions elliptiques de $\overline{\mathrm m}$ périodes, qui ont mêmes zéros et même poles, avec le $\overline{\mathrm m}$ ordre de multiplicité, sont égales à un facteur constant près ; Elles sont égales à une constante additive près, si elles ont $\overline{\mathrm m}$ périodes, et $\overline{\mathrm m}^{\mathrm s}$ poles, avec $\overline{\mathrm m}^{\mathrm s}$ parties principales.

\uline{Théorème $\overline{\mathrm{IV}}$} la somme des zéros d'une fonction \deleted{\unsure{holo|toute}} elliptique ds un parallélogramme de périodes égale la somme de ses poles, à une période près.

\uline{Théorème $\overline{\mathrm V}$}. \unsure{Ent|Entre} deux fonctions elliptiques de $\overline{\mathrm m}$ période Existe une relation algébrique

% archive sheet number: 60
